% ==============================================================================
%  Variedades, Fibración de Milnor y Temperatura en Hamiltonianos de Go
%  Análisis de 10 polinomios a través del diagrama de Hasse
%  Overleaf / pdfLaTeX compatible
% ==============================================================================
\documentclass[11pt,a4paper]{article}

% ── Codificación y fuente ──────────────────────────────────────────────────────
\usepackage[T1]{fontenc}
\usepackage[utf8]{inputenc}
\usepackage[spanish,es-noquoting]{babel}
\usepackage[default,scale=0.95]{sourcesanspro}
\usepackage[scaled=0.82]{inconsolata}
\usepackage{microtype}

% ── Geometría ─────────────────────────────────────────────────────────────────
\usepackage[top=2.2cm,bottom=2.4cm,left=2.4cm,right=2.4cm,
            headheight=14pt]{geometry}
\usepackage{parskip}
\setlength{\parskip}{5pt}

% ── Matemáticas ───────────────────────────────────────────────────────────────
\usepackage{amsmath,amssymb,amsthm,mathtools}

% ── Tablas ───────────────────────────────────────────────────────────────────
\usepackage{booktabs,tabularx,colortbl,array,multirow,longtable}
\usepackage{adjustbox}

% ── Color y gráficos ─────────────────────────────────────────────────────────
\usepackage[dvipsnames,table,x11names]{xcolor}
\usepackage{tikz}
\usetikzlibrary{positioning,arrows.meta,shapes,calc,decorations.pathmorphing,
                through,intersections,patterns}
\usepackage{tcolorbox}
\tcbuselibrary{skins,breakable,theorems}

% ── Tipografía ────────────────────────────────────────────────────────────────
\usepackage{titlesec}
\usepackage{fancyhdr}
\usepackage{enumitem}
\usepackage{hyperref}

% ── Paleta de colores ─────────────────────────────────────────────────────────
\definecolor{cNavy}{HTML}{1A2A4A}
\definecolor{cAccent}{HTML}{2E6DA4}
\definecolor{cGold}{HTML}{A07010}
\definecolor{cBg}{HTML}{F4F6FA}
\definecolor{cGrey}{HTML}{888888}
\definecolor{cLBlue}{HTML}{D6E8FA}
\definecolor{cGreen}{HTML}{1A7A3C}
\definecolor{cRed}{HTML}{8B1A1A}
\definecolor{cFront1}{HTML}{1A5C30}
\definecolor{cFront2}{HTML}{2E6DA4}
\definecolor{cFrontMid}{HTML}{7B5EA7}
\definecolor{cFrontLow}{HTML}{8B4513}
\definecolor{cFrontRef}{HTML}{444444}

\hypersetup{colorlinks=true,linkcolor=cAccent,urlcolor=cAccent,
            pdftitle={Variedades y Fibración de Milnor en Go}}

% ── Cajas tcolorbox ──────────────────────────────────────────────────────────
\tcbset{
  defbox/.style={enhanced,colback=cBg,colframe=cNavy,
    left=8pt,right=8pt,top=6pt,bottom=6pt,breakable,
    fonttitle=\bfseries\color{white},
    title style={fill=cNavy}},
  theorembox/.style={enhanced,colback=yellow!6,colframe=cGold,
    left=8pt,right=8pt,top=6pt,bottom=6pt,breakable,
    fonttitle=\bfseries\color{white},
    title style={fill=cGold}},
  hambox/.style={enhanced,colback=cLBlue!35,colframe=cAccent,
    left=8pt,right=8pt,top=5pt,bottom=5pt,breakable,
    fonttitle=\bfseries\color{cNavy},
    title style={fill=cLBlue,color=cNavy}},
  codebox/.style={enhanced,colback=cBg,colframe=cAccent!60,
    left=6pt,right=6pt,top=3pt,bottom=3pt,
    fontupper=\small\ttfamily,breakable},
  notebox/.style={enhanced,colback=yellow!10,colframe=cGold!60,
    left=6pt,right=6pt,top=3pt,bottom=3pt,breakable,
    fonttitle=\small\bfseries\color{cGold},
    title style={fill=cGold!20}}
}

% ── Secciones ─────────────────────────────────────────────────────────────────
\titleformat{\section}[block]{\large\bfseries\color{cNavy}}{%
  \color{cAccent}\S\,\thesection\quad}{0pt}{}%
  [\vspace{2pt}\color{cAccent}\rule{\linewidth}{1pt}\vspace{-4pt}]
\titleformat{\subsection}[hang]{\normalsize\bfseries\color{cNavy}}{%
  \color{cAccent}\thesubsection\quad}{0pt}{}
\titleformat{\subsubsection}[hang]{\small\bfseries\itshape\color{cNavy}}{}{0pt}{}

% ── Cabecera/pie ─────────────────────────────────────────────────────────────
\pagestyle{fancy}\fancyhf{}
\fancyhead[L]{\small\color{cGrey}Variedades, Fibración de Milnor y Temperatura --- Go}
\fancyhead[R]{\small\color{cGrey}\thepage}
\renewcommand{\headrule}{\color{cAccent}\hrule width\textwidth height 0.4pt}

% ── Operadores matemáticos ───────────────────────────────────────────────────
\DeclareMathOperator{\Hess}{Hess}
\DeclareMathOperator{\grad}{\nabla}
\newcommand{\R}{\mathbb{R}}
\newcommand{\Z}{\mathbb{Z}}
\newcommand{\C}{\mathbb{C}}
\newcommand{\crit}[1]{c^*_{#1}}
\newcommand{\Gamma}{\Gamma}
\newcommand{\dE}{\Delta E}
\newcommand{\Zpart}{Z(T)}
\newcommand{\pij}{p_{(s,s')}(T)}
\newcommand{\hbox}[1]{\texttt{#1}}

% ── Convenios de notación para pares de Go ──────────────────────────────────
\newcommand{\BB}{(\text{N},\text{N})}
\newcommand{\BV}{(\text{N},\varnothing)}
\newcommand{\BW}{(\text{N},\text{B})}
\newcommand{\WB}{(\text{B},\text{N})}
\newcommand{\WV}{(\text{B},\varnothing)}
\newcommand{\WW}{(\text{B},\text{B})}

% ══════════════════════════════════════════════════════════════════════════════
\begin{document}

% Barra de título
\begin{tikzpicture}[remember picture,overlay]
  \fill[cNavy] (current page.north west)
    rectangle ([yshift=-2.5cm]current page.north east);
  \node[anchor=west,text=white,font=\Large\bfseries] at
    ([xshift=1.8cm,yshift=-0.85cm]current page.north west)
    {Variedades, Fibración de Milnor y Temperatura};
  \node[anchor=west,text=cLBlue,font=\normalsize] at
    ([xshift=1.8cm,yshift=-1.55cm]current page.north west)
    {Análisis de 10 Hamiltonianos a través del Diagrama de Hasse de Go};
  \node[anchor=east,text=cGrey!70!white,font=\small] at
    ([xshift=-1.8cm,yshift=-2.15cm]current page.north east)
    {Experimento 06 $\cdot$ 2026-07-15};
\end{tikzpicture}
\vspace{2.6cm}

{\small\color{cGrey}\itshape
Catálogo: 296 candidatos evaluados $\cdot$
Plantilla: \texttt{cubic\_mixed} $\cdot$
Criterios: TDA ($H_1$), robustez, $\Delta E$, nodos $A_1$ $\cdot$
Diagrama de Hasse: 148 frentes Pareto
}

\tableofcontents
\newpage

% ==============================================================================
\section{Introducción: Go como sistema de espines}
% ==============================================================================

En el modelo de espín de Go, cada intersección del tablero lleva un valor
\[
  s \in \{-1,\,0,\,+1\}
  \quad\text{(negro, vacío, blanco)}
\]
y la energía de una configuración se calcula sumando sobre pares adyacentes:
\[
  E = \sum_{\langle i,j\rangle} H(s_i,\, s_j).
\]
Un \textbf{Hamiltoniano de par} $H:\{-1,0,+1\}^2 \to \R$ asigna un peso
energético a cada tipo de vecindad. La familia de todos los Hamiltonianos
posibles---no equivalentes sobre el dominio discreto---es exactamente la
familia de polinomios cúbicos: sobre $\{-1,0,+1\}$, la identidad $x^3=x$
reduce cualquier grado mayor a uno de los 7 monomios
\[
  x,\; y,\; x^2,\; xy,\; y^2,\; x^2y,\; xy^2.
\]
Esta es la \emph{plantilla \texttt{cubic\_mixed}}:
\[
  H(x,y) = a_1 x + a_2 y + b_{11}x^2 + b_{12}xy + b_{22}y^2
            + c_{112}x^2 y + c_{122}xy^2.
\]

En este informe analizamos \textbf{10 polinomios representativos} tomados de
distintas posiciones del diagrama de Hasse de orden parcial Pareto:
dos referencias históricas, dos candidatos del \textit{Frente~1}, y
candidatos de los Frentes 2, 3, 9, 25, 70 y 140.
Para cada uno describimos:
\begin{itemize}[noitemsep]
  \item la variedad $\Gamma(H)$ y sus zonas topológicas,
  \item los puntos críticos, nodos $A_1$ y cambios de topología de fibra,
  \item la construcción de la fibración de Milnor sobre $\Gamma(H)$,
  \item cómo mide la temperatura estratégica en Go.
\end{itemize}

% ==============================================================================
\section{Marco matemático}
% ==============================================================================

\subsection{La variedad \texorpdfstring{$\Gamma(H)$}{Gamma(H)} y sus fibras}

\begin{tcolorbox}[defbox,title={Definición: Variedad de un Hamiltoniano de Go}]
Dado $H:\R^2\to\R$, su \textbf{variedad asociada} es la superficie gráfica
\[
  \Gamma(H) \;=\; \{(x,\,y,\,z)\in\R^3 \mid z = H(x,y)\}.
\]
Es siempre una superficie suave (es un grafo sobre $\R^2$).
Para $c\in\R$ fijo, la \textbf{fibra de nivel} es
\[
  F_c \;=\; H^{-1}(c) \;=\; \{(x,y)\in\R^2 \mid H(x,y)=c\}\;\subset\;\R^2.
\]
Un valor $c$ es \textbf{regular} si $\nabla H(x,y)\neq 0$ para todo
$(x,y)\in F_c$; de lo contrario es un \textbf{valor crítico} $c^*$.
\end{tcolorbox}

La superficie $\Gamma(H)$ se divide en \emph{zonas} determinadas por los
valores críticos $c^*_1 < c^*_2 < \cdots < c^*_k$:
\begin{itemize}[noitemsep]
  \item \textbf{Zona superior}: $c > c^*_k$ --- fibra con una o más ramas no acotadas.
  \item \textbf{Zonas intermedias}: entre valores críticos consecutivos --- topología fija.
  \item \textbf{Zona inferior}: $c < c^*_1$ --- topología complementaria.
  \item \textbf{Frontera de zona} ($c = c^*$): fibra singular con nodos.
\end{itemize}

El \textbf{dominio Go} de $H$ son los 6 pares válidos (con al menos
una piedra):
$\{-1,+1\}\times\{-1,0,+1\}$,
que determinan los 6 valores energéticos
$E_{ss'} = H(s,s')$ sobre los cuales opera la dinámica estratégica.
Denotamos negro/blanco/vacío como N/B/$\varnothing$.

\subsection{Puntos críticos y nodos \texorpdfstring{$A_1$}{A1}}

\begin{tcolorbox}[defbox,title={Clasificación de puntos críticos}]
Un punto $(x_0,y_0)$ con $\nabla H(x_0,y_0)=0$ se clasifica por el determinante
del Hessiano:
\[
  \det\Hess H = H_{xx}H_{yy} - H_{xy}^2.
\]
\begin{itemize}[noitemsep]
  \item $\det\Hess H > 0$, $H_{xx}>0$: \textbf{mínimo local}.
  \item $\det\Hess H > 0$, $H_{xx}<0$: \textbf{máximo local}.
  \item $\det\Hess H < 0$: \textbf{punto de silla} o \textbf{nodo $A_1$}
        (singularidad aislada no degenerada). Es la singularidad genérica de un polinomio cúbico.
\end{itemize}
El \textbf{número de Milnor} $\mu$ es el número total de nodos $A_1$.
Para un cúbico con singularidades aisladas: $\mu \leq (d-1)^2 = 4$.
\end{tcolorbox}

En un nodo $A_1$ con valor crítico $c^*$, la fibra $F_{c^*}$ tiene una
auto-intersección transversal (figura en $\infty$ o \textit{figure-eight}),
que se denomina \textbf{ciclo de Van\-ishing}: un ciclo de $F_c$ (para
$c\neq c^*$ cercano) que se colapsa a un punto cuando $c\to c^*$.

\subsection{Fibración de Milnor: construcción}

\begin{tcolorbox}[theorembox,
  title={Fibración de Milnor para polinomios reales en $\R^2$}]
Sea $H:\R^2\to\R$ un polinomio con puntos críticos aislados
$p_1,\ldots,p_k$ de valores $c^*_1\leq\cdots\leq c^*_k$.

La \textbf{fibración de Milnor} de $H$ es la familia de fibras
\[
  \{F_c\}_{c\in\R\setminus\{c^*_1,\ldots,c^*_k\}}
\]
vista como fibrado localmente trivial sobre $\R\setminus\{c^*_i\}$.

Propiedades clave:
\begin{enumerate}[noitemsep]
  \item Para $c$ regular, $F_c$ es una curva suave de grado $d=3$ en $\R^2$.
  \item El género de la fibra genérica (sobre $\C$) es
        $g = \tfrac{(d-1)(d-2)}{2} = 1$ \textbf{(toro)}.
  \item Al cruzar un valor crítico $c^*_j$ (nodo $A_1$), la característica
        de Euler cambia: $\chi(F_c) \to \chi(F_c) \pm 1$.
  \item En términos topológicos: una componente oval se contrae a un punto
        y desaparece, o una rama se auto-intersecta y se divide.
  \item El número de componentes conexas de $F_c$ puede variar entre 1 y 3
        para un cúbico genérico.
\end{enumerate}
\end{tcolorbox}

\subsubsection{Interpretación topológica para Go}

Para un Hamiltoniano de Go cúbico, el diagrama de Hasse clasifica
$\mu$ (número de nodos $A_1$) como uno de los cuatro criterios Pareto.
Cada nodo $A_1$ representa un \textbf{cambio de régimen estratégico}:
la configuración de pares de piedras que el modelo distingue
pasa por una transición topológica en la variedad.

\begin{itemize}[noitemsep]
  \item $\mu = 0$: H sin puntos silla $\Rightarrow$ estrategia monotónica
        (preferencia global sin inversiones).
  \item $\mu = 1$: un cambio de régimen $\Rightarrow$ un umbral estratégico.
  \item $\mu = 2$: dos umbrales, el máximo posible para los candidatos
        del catálogo.
\end{itemize}

\subsection{Temperatura y función de partición}

\begin{tcolorbox}[defbox,title={Temperatura en el modelo de espín de Go}]
Fijado $H$ y una temperatura $T>0$, la \textbf{función de partición} sobre
los 6 pares válidos es:
\[
  Z(T) = \sum_{(s,s')\in\{-1,+1\}\times\{-1,0,+1\}}
         e^{-H(s,s')/T}.
\]
La \textbf{probabilidad de ocupación} del par $(s,s')$ es:
\[
  p_{(s,s')}(T) = \frac{e^{-H(s,s')/T}}{Z(T)}.
\]
La \textbf{energía media} y la \textbf{capacidad calorífica} son:
\[
  \langle H\rangle(T) = \sum_{(s,s')} H(s,s')\,p_{(s,s')}(T), \qquad
  C(T) = \frac{\text{Var}(H)}{T^2} = \frac{\langle H^2\rangle - \langle H\rangle^2}{T^2}.
\]
El rango energético $\dE = \max H - \min H$ sobre los 6 pares fija la
\textbf{escala de temperatura estratégica}: para $T\ll\dE$ el modelo juega
determinísticamente (elige el par de menor energía); para $T\gg\dE$ juega
aleatoriamente.
\end{tcolorbox}

La persistencia TDA ($H_1$-máxima) mide la \textbf{complejidad topológica
térmica}: cuántas ``burbujas'' en el paisaje de energía persisten a través
de una gama de temperaturas antes de colapsar.

\newpage
% ==============================================================================
\section{Los 10 Hamiltonianos: análisis individual}
% ==============================================================================

Presentamos los 10 polinomios en orden de posición en el diagrama de Hasse,
desde los dos modelos de referencia histórica hasta el candidato del Frente
140 (dominado por casi todos los demás).

% ─────────────────────────────────────────────────────────────────────────────
\subsection{\texorpdfstring{$H_{\text{Alv}}$}{H\_Alv} --- Modelo Alvarado (Referencia, F-REF)}
% ─────────────────────────────────────────────────────────────────────────────

\begin{tcolorbox}[hambox,title={$H_{\text{Alv}}$ (Alvarado, Atomic-Go) $\cdot$ Frente: Referencia}]
\[
  H_{\text{Alv}}(x,y) = xy
\]
\textbf{Plantilla}: \texttt{quadratic} $\quad$
\textbf{Nodos $A_1$}: 1 $\quad$
\textbf{Robustez}: 0.0 $\quad$
\textbf{Score}: 0.090
\end{tcolorbox}

\subsubsection*{Variedad $\Gamma(H_{\text{Alv}})$ y sus zonas}

La superficie $z = xy$ es una \textbf{silla de montar} (paraboloide hiperbólico):
\[
  z = xy = \tfrac{1}{2}\bigl[(x+y)^2/2 - (x-y)^2/2\bigr].
\]
Las curvas de nivel $xy = c$ son hipérbolas (para $c\neq 0$) o las dos rectas
$x=0$ e $y=0$ (para $c=0$). El único punto crítico es el origen:
\[
  \nabla H_{\text{Alv}} = (y, x) = (0,0) \quad\Longrightarrow\quad (x,y)=(0,0).
\]
Con $H_{xx}=0$, $H_{yy}=0$, $H_{xy}=1$: $\det\Hess = -1 < 0$ $\Rightarrow$ \textbf{nodo $A_1$}, valor crítico $c^*=0$.

\medskip\noindent
\textbf{Zonas de $\Gamma(H_{\text{Alv}})$}:
\begin{itemize}[noitemsep]
  \item \textbf{Zona positiva} ($c>0$): las hipérbolas $xy=c>0$ tienen dos ramas
        en los cuadrantes I y III (mismos signos: NN y BB).
  \item \textbf{Frontera} ($c=0$): los ejes coordenados; el vacío ($s=0$)
        da $H=0$ para cualquier vecino.
  \item \textbf{Zona negativa} ($c<0$): hipérbolas $xy=c<0$ en los cuadrantes II y IV
        (signos opuestos: NB y BN).
\end{itemize}

\subsubsection*{Fibración de Milnor}

Para $c\neq 0$, la fibra $F_c = \{xy=c\}$ es una hipérbola suave: dos ramas
no acotadas, cada una homeomorfa a $\R$ (topología: dos copias de $\R$).
En $c=0$, la fibra degenera en dos rectas que se cruzan: la figura en `X'
(singularidad tipo $A_1$). El ciclo de Vanishing es el ``bucle'' que
rodea el origen y colapsa a la X al llegar a $c^*=0$.

\subsubsection*{Energías en pares de Go}

\begin{center}
\begin{tabular}{cccccc}
  \toprule
  Par & $(-1,-1)$ & $(-1,0)$ & $(-1,+1)$ & $(+1,-1)$ & $(+1,0)$ & $(+1,+1)$ \\
  \midrule
  Notación & NN & N$\varnothing$ & NB & BN & B$\varnothing$ & BB \\
  $H_{\text{Alv}}$ & $+1$ & $0$ & $-1$ & $-1$ & $0$ & $+1$ \\
  \bottomrule
\end{tabular}
\end{center}

$\dE=2$. El modelo premia la vecindad entre \emph{piedras del mismo color}
(energía $-1$: favorecida) y penaliza la vecindad homogénea de mismo signo
(¡wait: $H=+1$ para NN y BB es la energía \emph{más alta}!).

\begin{tcolorbox}[notebox,title={Lectura estratégica}]
$H_{\text{Alv}}$ es \textbf{simétrico}: $H(-x,-y)=H(x,y)$ (función par).
No distingue negro de blanco.
Minimiza la energía colocando piedras de colores opuestos como vecinas.
Representa el \textbf{Ising antiferromagnético}: el vacío es invisible ($H(s,0)=0$).
Temperatura característica: $T_{\text{eff}} \approx \dE/\ln 2 \approx 1.44$.
\end{tcolorbox}

% ─────────────────────────────────────────────────────────────────────────────
\subsection{\texorpdfstring{$H_{\text{M1}}$}{H\_M1} --- Modelo Mercado-Jiménez (Referencia, F-REF)}
% ─────────────────────────────────────────────────────────────────────────────

\begin{tcolorbox}[hambox,title={$H_{\text{M1}}$ (Mercado \& Jiménez, 2024) $\cdot$ Frente: Referencia}]
\[
  H_{\text{M1}}(x,y) = x + 2y - x^2y - xy^2
\]
\textbf{Plantilla}: \texttt{sparse\_cubic} $\quad$
\textbf{Nodos $A_1$}: 2 $\quad$
\textbf{Robustez}: 0.333 $\quad$
\textbf{Score}: 0.190
\end{tcolorbox}

\subsubsection*{Variedad $\Gamma(H_{\text{M1}})$ y sus zonas}

La superficie $z = x + 2y - x^2y - xy^2 = x(1-xy) + y(2-x^2)$
exhibe un crecimiento asintótico cúbico dominado por $-x^2y - xy^2$.

El gradiente es:
\[
  \frac{\partial H}{\partial x} = 1 - 2xy - y^2, \qquad
  \frac{\partial H}{\partial y} = 2 - x^2 - 2xy.
\]
Igualando a cero se obtienen dos puntos críticos simétricos:
$(x_0,y_0)\approx (\pm 1.075,\, \pm 0.393)$
con valores críticos $c^* = \mp 1.2408$.
La función es impar: $H_{\text{M1}}(-x,-y) = -H_{\text{M1}}(x,y)$,
lo que garantiza la simetría $c^*_2 = -c^*_1 = 1.2408$.

\medskip\noindent
\textbf{Zonas de $\Gamma(H_{\text{M1}})$}:

\begin{itemize}[noitemsep]
  \item $c > 1.2408$: fibra desconectada en dos ramas separadas.
  \item $c = 1.2408$: primera singularidad $A_1$; nodo en $(1.075, 0.393)$.
  \item $-1.2408 < c < 1.2408$: fibra reconectada; topología de ciclo.
  \item $c = -1.2408$: segunda singularidad $A_1$; nodo en $(-1.075,-0.393)$.
  \item $c < -1.2408$: misma topología que la zona superior (por simetría impar).
\end{itemize}

\subsubsection*{Fibración de Milnor}

Hay exactamente $\mu=2$ nodos $A_1$, separados por $\Delta c^* = 2.4816$.
La fibra genérica de un cúbico tiene género $g=1$ (toro) sobre $\C$.
En la región $|c| < 1.2408$, la fibra real $F_c$ contiene un oval compacto
(ciclo no trivial): es la manifestación real del asa toroidal.
Al cruzar $c^*=\pm 1.2408$, este oval se colapsa a un punto (ciclo de Vanishing)
y la topología cambia de ``oval + rama'' a ``solo rama''.

\subsubsection*{Energías en pares de Go}

\begin{center}
\begin{tabular}{cccccc}
  \toprule
  Par & NN & N$\varnothing$ & NB & BN & B$\varnothing$ & BB \\
  \midrule
  $H_{\text{M1}}$ & $-1$ & $-1$ & $+1$ & $-1$ & $+1$ & $+1$ \\
  \bottomrule
\end{tabular}
\end{center}

$\dE=2$. Solo 2 niveles de energía: $\{-1, +1\}$.
El modelo distingue \emph{pares que contienen una piedra blanca como segundo elemento}
(energía $+1$: penalizados) de todos los demás (energía $-1$: favorecidos).

\begin{tcolorbox}[notebox,title={Lectura estratégica}]
$H_{\text{M1}}$ es \textbf{asimétrico}: negro ($s=-1$) como piedra activa
tiene ventaja estructural sobre blanco.
El vacío como segundo vecino da $H=-1$ si el primero es negro,
$H=+1$ si es blanco: el \textbf{vacío es activo} (no invisible).
Temperatura característica: $T_{\text{eff}}\approx 1.44$.
\end{tcolorbox}

% ─────────────────────────────────────────────────────────────────────────────
\subsection{\texorpdfstring{$H_{\text{F1a}}$}{H\_F1a} --- Candidato H\_0045 (Frente 1)}
% ─────────────────────────────────────────────────────────────────────────────

\begin{tcolorbox}[hambox,title={$H_{\text{F1a}}$ (H\_0045, \textcolor{cFront1}{Frente 1}) $\cdot$
  $H_1$-max = 0.1664 $\cdot$ Rob = 1.000}]
\begin{align*}
  H_{\text{F1a}}(x,y) &=
    -0.7583\,x + 0.1425\,y
    - 2.3900\,x^2 + 2.0008\,xy - 2.6882\,y^2 \\
  &\quad + 0.8497\,x^2y - 0.8018\,xy^2
\end{align*}
\textbf{Nodos $A_1$}: 2 $\quad$
\textbf{$\dE$}: 8.00 $\quad$
\textbf{Score}: 0.454
\end{tcolorbox}

\subsubsection*{Variedad y zonas}

El término dominante en escala grande es $-2.390\,x^2 - 2.688\,y^2$
(paraboloide elíptico negativo), modulado por los cúbicos $+0.850\,x^2y$
y $-0.802\,xy^2$.
Cerca del origen, la superficie desciende en todas las direcciones.
Los dos nodos $A_1$ se ubican \emph{lejos del dominio Go}:
\[
  p_1 \approx (-0.171,\,-0.034),\quad c^*_1 \approx 0.063\quad (\text{máximo local})
\]
\[
  p_2 \approx (6.288,\,2.996),\quad c^*_2 \approx -29.88\quad (A_1\text{-nodo})
\]
\[
  p_3 \approx (-2.957,\,2.613),\quad c^*_3 \approx -16.50\quad (A_1\text{-nodo})
\]

\textbf{Zonas de $\Gamma(H_{\text{F1a}})$}:

\begin{itemize}[noitemsep]
  \item $c > 0.063$: la fibra está por encima del máximo local; vacía en el
        dominio Go central.
  \item $-16.50 < c < 0.063$: fibra con topología de \textbf{oval cerrado}
        que contiene los 6 pares Go (todos con $-9.63 \leq E \leq -1.63$).
  \item $-29.88 < c < -16.50$: una rama adicional aparece por el nodo en
        $(6.29, 3.00)$; el oval Go sigue intacto.
  \item $c < -29.88$: segunda transición topológica.
\end{itemize}

\subsubsection*{Milnor y ciclos de Vanishing}

Los dos nodos $A_1$ activos ($\mu=2$) generan \textbf{dos ciclos de Vanishing}.
El ciclo de Vanishing del nodo en $(-2.957, 2.613)$ (con $c^*=-16.50$)
se encuentra dentro del rango de energía de los pares Go, ya que
$E_{\min}=-9.63 > -16.50$. Esto significa que la fibra ``vista'' por
los pares Go incluye el ciclo que va a colapsar al cruzar $-16.50$:
hay \textbf{complejidad topológica accesible} para el sistema Go.

\subsubsection*{Energías en pares de Go}

\begin{center}
\begin{tabular}{ccccccc}
  \toprule
  Par & NN & N$\varnothing$ & NB & BN & B$\varnothing$ & BB & $\dE$ \\
  \midrule
  $H_{\text{F1a}}$ & $-2.51$ & $-1.63$ & $-4.53$ & $-9.63$ & $-3.15$ & $-3.65$ & 8.00 \\
  \bottomrule
\end{tabular}
\end{center}

Todos los pares tienen energía negativa. El par BN ($s_i=+1, s_j=-1$)
es el más favorecido ($E=-9.63$): \textbf{piedra blanca con vecino negro}
es la configuración de mínima energía.

\begin{tcolorbox}[notebox,title={Lectura estratégica y temperatura}]
La asimetría BN $\ll$ NB (razón $\approx 2.1$) introduce una
\textbf{direccionalidad}: colocar blanco junto a negro vale mucho más que
negro junto a blanco. La escala de temperatura es $T\approx\dE/\ln 6\approx 4.5$.
La persistencia $H_1=0.166$ indica que la burbuja topológica Go dura $\approx 0.166$
unidades de temperatura antes de colapsar: buen discriminador térmico.
\end{tcolorbox}

% ─────────────────────────────────────────────────────────────────────────────
\subsection{\texorpdfstring{$H_{\text{F1b}}$}{H\_F1b} --- Candidato H\_0042 (Frente 1)}
% ─────────────────────────────────────────────────────────────────────────────

\begin{tcolorbox}[hambox,title={$H_{\text{F1b}}$ (H\_0042, \textcolor{cFront1}{Frente 1}) $\cdot$
  $H_1$-max = 0.1204 $\cdot$ Rob = 1.000}]
\begin{align*}
  H_{\text{F1b}}(x,y) &=
    -1.2917\,x - 2.0783\,y
    - 2.3071\,x^2 - 2.8731\,xy - 2.6676\,y^2 \\
  &\quad - 0.6507\,x^2y - 0.8932\,xy^2
\end{align*}
\textbf{Nodos $A_1$}: 2 $\quad$
\textbf{$\dE$}: 11.75 $\quad$
\textbf{Score}: 0.417
\end{tcolorbox}

\subsubsection*{Variedad y zonas}

Todos los coeficientes son negativos. La superficie
$z = H_{\text{F1b}}(x,y)$ desciende en todas las direcciones desde un
pequeño máximo local en $(-0.095, -0.351)$ con $c^*_{\max}\approx 0.419$.
Los dos nodos $A_1$ están en regiones \emph{muy alejadas} del dominio Go:
\[
  p_1\approx(-3.089,\,-3.221),\quad c^*_1=-18.97;\qquad
  p_2\approx(-2.880,\,+4.187),\quad c^*_2=-13.74.
\]

\textbf{Zonas}:
\begin{itemize}[noitemsep]
  \item $c > 0.419$: fibra vacía (no hay $(x,y)$ con $H=c>0.419$
        en la región del dominio Go).
  \item $-12.76 \leq c \leq 0.419$: los 6 pares Go están en este rango.
        La fibra es un oval compacto que encierra el dominio Go.
  \item $-13.74 < c < -12.76$: primer nodo $A_1$ aún no alcanzado.
  \item $c = -13.74$: \textbf{primer nodo $A_1$}: la fibra adquiere una
        auto-intersección justo por debajo del par BB ($E=-12.76$).
  \item $c = -18.97$: \textbf{segundo nodo $A_1$}.
\end{itemize}

\subsubsection*{Milnor: complejidad accesible}

El par BB con $E=-12.76$ está \emph{entre} $c^*_2=-13.74$ y
$c^*_3=-18.97$: la fibra al nivel del par más favorecido está
precisamente en la zona de transición topológica.
Esto explica la persistencia $H_1=0.120$: hay un ciclo topológico
cuya ``vida'' corresponde a la separación energética $12.76-0$ del máximo.

\subsubsection*{Energías}

\begin{center}
\begin{tabular}{ccccccc}
  \toprule
  Par & NN & N$\varnothing$ & NB & BN & B$\varnothing$ & BB & $\dE$ \\
  \midrule
  $H_{\text{F1b}}$ & $-2.93$ & $-1.02$ & $-2.65$ & $-1.56$ & $-3.60$ & $-12.76$ & 11.75 \\
  \bottomrule
\end{tabular}
\end{center}

BB es el mínimo absoluto con una separación notable: $-12.76$ vs.\ los demás
($-1.0$ a $-3.6$). El modelo concentra toda la preferencia en la
\textbf{vecindad blanco-blanco}.

\begin{tcolorbox}[notebox,title={Lectura estratégica y temperatura}]
Un modelo con esta distribución favorece abrumadoramente las formaciones
compactas de piedras blancas. A temperatura $T\approx 2$, la probabilidad de
elegir el par BB supera el 99\%. La escala $\dE=11.75$ da excelente
resolución térmica: el paso de juego aleatorio a determinístico ocurre
sobre una ventana de temperatura muy precisa.
\end{tcolorbox}

% ─────────────────────────────────────────────────────────────────────────────
\subsection{\texorpdfstring{$H_{\text{F2}}$}{H\_F2} --- Candidato H\_0022 (Frente 2)}
% ─────────────────────────────────────────────────────────────────────────────

\begin{tcolorbox}[hambox,title={$H_{\text{F2}}$ (H\_0022, \textcolor{cFront2}{Frente 2}) $\cdot$
  $H_1$-max = 0.1798 $\cdot$ Rob = 1.000 $\cdot$ \textbf{Mayor $H_1$ del grupo}}]
\begin{align*}
  H_{\text{F2}}(x,y) &=
    -1.0208\,x - 2.1329\,y
    - 2.3796\,x^2 + 0.5259\,xy - 1.9764\,y^2 \\
  &\quad + 0.8502\,x^2y + 0.1621\,xy^2
\end{align*}
\textbf{Nodos $A_1$}: \textbf{1} $\quad$
\textbf{$\dE$}: 4.61 $\quad$
\textbf{Score}: 0.458
\end{tcolorbox}

\subsubsection*{Variedad y zonas}

Con solo $\mu=1$ nodo $A_1$, este es el Hamiltoniano topológicamente más
\emph{simple} entre los candidatos del Frente~2. El único punto silla está en:
\[
  p_1\approx(3.072,\,2.539),\quad c^*_1=-16.06.
\]
Hay además un \textbf{máximo local} en $(-0.222,\,-0.549)$ con
$c^*_{\max}=+0.715$.

\textbf{Zonas}:
\begin{itemize}[noitemsep]
  \item $c > 0.715$: fibra vacía en el dominio Go.
  \item $-5.97 \leq c \leq 0.715$: los 6 pares Go están en esta zona.
        Un solo oval compacto con una sola transición futura.
  \item $c = -16.06$: único nodo $A_1$, lejos del dominio Go.
  \item $c < -16.06$: cambio topológico; una rama adicional aparece.
\end{itemize}

\subsubsection*{Por qué $H_1=0.1798$ es el más alto}

La persistencia TDA mide el tiempo de vida de los ``loops'' en la
filtración por subniveles de $H$ sobre la cuadrícula discreta.
Con $\mu=1$ y el nodo $A_1$ en $c^*=-16.06$ (bien separado de los pares Go),
el único cambio de topología de la fibra dentro del rango de energía Go
es suave y prolongado: el oval del dominio Go vive desde $c\approx -5.97$
hasta $c\approx 0.715$, una longitud de $\approx 6.69$ niveles.
El pico de persistencia $H_1=0.180$ refleja este ciclo de larga duración.

\subsubsection*{Energías}

\begin{center}
\begin{tabular}{ccccccc}
  \toprule
  Par & NN & N$\varnothing$ & NB & BN & B$\varnothing$ & BB & $\dE$ \\
  \midrule
  $H_{\text{F2}}$ & $-1.69$ & $-1.36$ & $-5.31$ & $-4.46$ & $-3.40$ & $-5.97$ & 4.61 \\
  \bottomrule
\end{tabular}
\end{center}

Los pares mixtos (NB, BN, BB) son los más favorecidos.

% ─────────────────────────────────────────────────────────────────────────────
\subsection{\texorpdfstring{$H_{\text{F3}}$}{H\_F3} --- Candidato H\_0027 (Frente 3)}
% ─────────────────────────────────────────────────────────────────────────────

\begin{tcolorbox}[hambox,title={$H_{\text{F3}}$ (H\_0027, Frente 3) $\cdot$
  $H_1$-max = 0.0072 $\cdot$ Rob = 1.000}]
\begin{align*}
  H_{\text{F3}}(x,y) &=
    +2.4961\,x - 1.6187\,y
    - 2.7755\,x^2 + 0.3291\,xy - 0.7745\,y^2 \\
  &\quad + 0.6596\,x^2y + 0.6165\,xy^2
\end{align*}
\textbf{Nodos $A_1$}: 2 $\quad$
\textbf{$\dE$}: 8.14 $\quad$
\textbf{Score}: 0.392
\end{tcolorbox}

\subsubsection*{Variedad y zonas}

Los términos lineales $+2.496\,x - 1.619\,y$ inclinan fuertemente la
superficie: positivo en la dirección $x$ y negativo en $y$.
Los puntos críticos están en posiciones moderadas:
\[
  p_1 \approx (0.431,-1.331),\quad c^*_{\max}=+1.462;
\]
\[
  p_2 \approx (1.212,-4.615),\quad c^*=-0.474\;(A_1);
\]
\[
  p_3 \approx (1.298,+1.566),\quad c^*=-1.499\;(A_1).
\]
El rango de valores críticos es pequeño: $[-1.499, 1.462]$.
Los pares Go tienen energías en $[-7.95, +0.19]$, por lo que los
valores críticos caen \emph{dentro} del rango de los pares Go.

\textbf{Zonas y su intersección con el dominio Go}:
\begin{itemize}[noitemsep]
  \item $c > 1.462$: solo el par BN ($E=+0.193$) y
        B$\varnothing$ ($E=-0.280$) están en esta región.
  \item $-0.474 < c < 1.462$: máximo local + primer nodo $A_1$ en este
        rango; todos los pares con $E > -0.474$.
  \item $-1.499 < c < -0.474$: segundo nodo $A_1$; los pares NN, N$\varnothing$
        y NB entran en esta zona.
  \item $c < -1.499$: zona estable, contiene el mínimo global NB ($E=-7.95$).
\end{itemize}

\subsubsection*{Milnor: dos transiciones dentro del dominio Go}

Este es el caso más interesante topológicamente para Go:
\textbf{ambos nodos $A_1$ tienen valores críticos dentro del rango de
pares válidos de Go}.
Esto significa que al barrer la temperatura de $c=-7.95$ a $c=+0.19$,
el sistema Go experimenta \emph{dos} cambios de régimen topológico,
en $c^*=-1.499$ y $c^*=-0.474$.

Sin embargo, la persistencia TDA es solo $H_1=0.007$: los ciclos son
efímeros. La separación entre valores críticos es solo
$\Delta c^* = |-0.474-(-1.499)|=1.025$, lo que implica que la
burbuja topológica intermedia es muy angosta.

\subsubsection*{Energías}

\begin{center}
\begin{tabular}{ccccccc}
  \toprule
  Par & NN & N$\varnothing$ & NB & BN & B$\varnothing$ & BB & $\dE$ \\
  \midrule
  $H_{\text{F3}}$ & $-5.37$ & $-5.27$ & $-7.95$ & $+0.19$ & $-0.28$ & $-1.07$ & 8.14 \\
  \bottomrule
\end{tabular}
\end{center}

El único par con energía \emph{positiva} es BN ($+0.19$):
el modelo penaliza levemente la vecindad ``piedra blanca con negro al lado''.

% ─────────────────────────────────────────────────────────────────────────────
\subsection{\texorpdfstring{$H_{\text{F9}}$}{H\_F9} --- Candidato H\_0271 (Frente 9)}
% ─────────────────────────────────────────────────────────────────────────────

\begin{tcolorbox}[hambox,title={$H_{\text{F9}}$ (H\_0271, \textcolor{cFrontMid}{Frente 9}) $\cdot$
  $H_1$-max = 0.000 $\cdot$ Rob = 0.000}]
\begin{align*}
  H_{\text{F9}}(x,y) &=
    +1.9970\,x - 1.1808\,y
    - 0.3611\,x^2 + 2.9137\,xy + 2.9103\,y^2 \\
  &\quad + 0.8122\,x^2y + 0.5969\,xy^2
\end{align*}
\textbf{Nodos $A_1$}: 2 $\quad$
\textbf{$\dE$}: 11.02 $\quad$
\textbf{Score}: 0.090
\end{tcolorbox}

\subsubsection*{Variedad y zonas}

El término $+2.910\,y^2$ hace que la superficie \emph{crezca} en la
dirección $y$, en contraste con los candidatos del Frente~1.
Los valores críticos son \emph{positivos}: $c^*_1=1.186$ y $c^*_2=14.25$.
Los pares Go tienen energías en $[-3.33, +7.69]$, que incluyen el primer
valor crítico $c^*_1=1.186$ pero no el segundo.

\textbf{Zonas relevantes para Go}:
\begin{itemize}[noitemsep]
  \item $c < 1.186$: los pares N$\varnothing$, NB (negativos) están aquí.
        Un nodo $A_1$ separa esta zona de la siguiente.
  \item $1.186 < c < 7.69$: los pares BN, B$\varnothing$, NN, BB (positivos)
        están aquí. El ciclo de Vanishing ya colapsó.
  \item $c > 14.25$: segundo nodo $A_1$, fuera del dominio Go.
\end{itemize}

\subsubsection*{Por qué Rob = 0 y $H_1$ = 0}

La robustez cero indica que \emph{ninguna} de las 12 perturbaciones
$\pm 5\%$ de los coeficientes mantiene $H_1 > 0.20$.
La persistencia $H_1=0$ significa que no hay ciclos de larga vida
en la filtración TDA: el paisaje energético no tiene ``cuencas''
topológicamente no triviales en el rango discreto de Go.
Este candidato fue superado por los del Frente~1 precisamente en
este criterio: el nodo $A_1$ en $c^*=1.186$ cae dentro del rango Go
pero la burbuja que genera es instantánea ($H_1=0$).

\subsubsection*{Energías}

\begin{center}
\begin{tabular}{ccccccc}
  \toprule
  Par & NN & N$\varnothing$ & NB & BN & B$\varnothing$ & BB & $\dE$ \\
  \midrule
  $H_{\text{F9}}$ & $+3.24$ & $-2.36$ & $-3.33$ & $+2.60$ & $+1.64$ & $+7.69$ & 11.02 \\
  \bottomrule
\end{tabular}
\end{center}

Los pares ``mismo color'' (NN=$+3.24$, BB=$+7.69$) son altamente penalizados.
La vecindad entre piedras del mismo color es muy costosa energéticamente.

% ─────────────────────────────────────────────────────────────────────────────
\subsection{\texorpdfstring{$H_{\text{F25}}$}{H\_F25} --- Candidato H\_0185 (Frente 25)}
% ─────────────────────────────────────────────────────────────────────────────

\begin{tcolorbox}[hambox,title={$H_{\text{F25}}$ (H\_0185, \textcolor{cFrontMid}{Frente 25}) $\cdot$
  $H_1$-max = 0.000 $\cdot$ Rob = 0.000}]
\begin{align*}
  H_{\text{F25}}(x,y) &=
    +2.8718\,x - 0.0082\,y
    - 0.0380\,x^2 + 2.8253\,xy - 0.7156\,y^2 \\
  &\quad - 0.2049\,x^2y + 0.1695\,xy^2
\end{align*}
\textbf{Nodos $A_1$}: 2 $\quad$
\textbf{$\dE$}: 11.73 $\quad$
\textbf{Score}: 0.090
\end{tcolorbox}

\subsubsection*{Variedad y zonas}

El término dominante lineal $+2.872\,x$ crea una fuerte inclinación
en la dirección $x$: la superficie sube rápidamente para $x>0$
y baja para $x<0$.
El coeficiente cuadrático dominante es $+2.825\,xy$:
la interacción cruzada directa (tipo Ising).

Los valores críticos son muy asimétricos: $c^*_1=-0.776$ (cercano)
y $c^*_2=33.01$ (muy lejano, en $(14.47, 0.59)$).

\textbf{Zonas}:
\begin{itemize}[noitemsep]
  \item $c < -0.776$: los pares NN, N$\varnothing$, NB (negativos) están aquí;
        primer nodo $A_1$ separa la zona negativa.
  \item $-0.776 < c < 4.90$: los pares BN ($-0.32$), B$\varnothing$ ($+2.83$),
        BB ($+4.90$) están aquí.
  \item $c > 33.01$: segundo nodo $A_1$, muy lejos del dominio Go.
\end{itemize}

El nodo $A_1$ en $c^*=-0.776$ cae entre los pares negativos y positivos,
marcando el \textbf{umbral entre colocar negro y colocar blanco}.

\subsubsection*{Energías}

\begin{center}
\begin{tabular}{ccccccc}
  \toprule
  Par & NN & N$\varnothing$ & NB & BN & B$\varnothing$ & BB & $\dE$ \\
  \midrule
  $H_{\text{F25}}$ & $-0.76$ & $-2.91$ & $-6.83$ & $-0.32$ & $+2.83$ & $+4.90$ & 11.73 \\
  \bottomrule
\end{tabular}
\end{center}

\begin{tcolorbox}[notebox,title={Lectura estratégica}]
Este modelo favorece fuertemente los pares donde el \emph{primer} elemento
es negro (NB=$-6.83$, N$\varnothing$=$-2.91$): negro con vecino blanco es
óptimo. Los pares con primer elemento blanco y segundo positivo
(B$\varnothing$, BB) son penalizados. Es un modelo con fuerte asimetría
de color: negro domina como piedra activa.
\end{tcolorbox}

% ─────────────────────────────────────────────────────────────────────────────
\subsection{\texorpdfstring{$H_{\text{F70}}$}{H\_F70} --- Candidato H\_0004 (Frente 70)}
% ─────────────────────────────────────────────────────────────────────────────

\begin{tcolorbox}[hambox,title={$H_{\text{F70}}$ (H\_0004, \textcolor{cFrontLow}{Frente 70}) $\cdot$
  $H_1$-max = 0.000 $\cdot$ Rob = 0.000}]
\begin{align*}
  H_{\text{F70}}(x,y) &=
    +1.6437\,x - 0.3667\,y
    + 2.1516\,x^2 + 1.1842\,xy - 2.4349\,y^2 \\
  &\quad + 0.9512\,x^2y + 0.5223\,xy^2
\end{align*}
\textbf{Nodos $A_1$}: \textbf{1} $\quad$
\textbf{$\dE$}: 6.84 $\quad$
\textbf{Score}: 0.090
\end{tcolorbox}

\subsubsection*{Variedad y zonas}

El término $+2.152\,x^2$ hace que la superficie \emph{suba} en ambas
direcciones de $x$ (coeficiente positivo), mientras $-2.435\,y^2$ la
hace bajar en $y$. El único nodo $A_1$ está en:
\[
  p_1\approx(-0.370,\,-0.128),\quad c^*=-0.270.
\]
Este punto está muy \emph{cerca del origen} y dentro del dominio relevante
de los pares Go. Con un solo valor crítico ($\mu=1$), la fibración
de Milnor tiene exactamente una transición.

\textbf{Zonas}:
\begin{itemize}[noitemsep]
  \item $c < -3.05$: los pares NB ($-3.05$) y NN ($-1.85$) están aquí;
        topología pre-transición.
  \item $c = -0.270$: único nodo $A_1$; el ciclo de Vanishing colapsa.
  \item $c > -0.270$: los pares B$\varnothing$ ($+3.80$), BB ($+3.65$),
        BN ($+0.11$) y N$\varnothing$ ($+0.51$) están en zona post-transición.
\end{itemize}

\subsubsection*{Energías}

\begin{center}
\begin{tabular}{ccccccc}
  \toprule
  Par & NN & N$\varnothing$ & NB & BN & B$\varnothing$ & BB & $\dE$ \\
  \midrule
  $H_{\text{F70}}$ & $-1.85$ & $+0.51$ & $-3.05$ & $+0.11$ & $+3.80$ & $+3.65$ & 6.84 \\
  \bottomrule
\end{tabular}
\end{center}

El nodo $A_1$ en $c^*=-0.270$ separa el grupo de pares
negativos/bajos ($\text{NN}=-1.85$, $\text{NB}=-3.05$) del grupo positivo
($\text{B}\varnothing=+3.80$, $\text{BB}=+3.65$). La transición topológica
coincide con el umbral estratégico ``negro vs.\ blanco como primer elemento''.

% ─────────────────────────────────────────────────────────────────────────────
\subsection{\texorpdfstring{$H_{\text{F140}}$}{H\_F140} --- Candidato H\_0013 (Frente 140)}
% ─────────────────────────────────────────────────────────────────────────────

\begin{tcolorbox}[hambox,title={$H_{\text{F140}}$ (H\_0013, \textcolor{cFrontLow}{Frente 140}) $\cdot$
  $H_1$-max = 0.000 $\cdot$ Rob = 0.000}]
\begin{align*}
  H_{\text{F140}}(x,y) &=
    +1.5900\,x + 0.8083\,y
    + 0.3215\,x^2 + 0.3552\,xy - 1.1763\,y^2 \\
  &\quad - 0.9384\,x^2y - 0.1266\,xy^2
\end{align*}
\textbf{Nodos $A_1$}: 2 $\quad$
\textbf{$\dE$}: 4.72 $\quad$
\textbf{Score}: 0.090
\end{tcolorbox}

\subsubsection*{Variedad y zonas}

Este candidato presenta un fenómeno notable: uno de sus nodos $A_1$ está
en una posición \emph{extremadamente alejada} del origen:
\[
  p_2\approx(-12.44,\,186.96),\quad c^*_2=-13876.85.
\]
El otro está en:
\[
  p_1\approx(-1.107,\,-0.355),\quad c^*_1=-1.236.
\]
El valor crítico $c^*_2=-13877$ es numéricamente gigantesco, producto de
los términos cúbicos $-0.938\,x^2y$ cuando $y=186.96$. Para todos los
propósitos Go, \textbf{solo el nodo cercano} con $c^*=-1.236$ es relevante.

\textbf{Zonas relevantes para Go}:
\begin{itemize}[noitemsep]
  \item $c < -2.80$: solo NB ($-2.80$) aquí.
  \item $-2.80 < c < -1.24$: NN, N$\varnothing$ y NB por encima de $-1.24$.
  \item $c = -1.236$: nodo $A_1$ cercano; transición topológica.
  \item $c > -1.236$: los pares BN ($+0.38$), BB ($+0.83$), B$\varnothing$ ($+1.91$).
\end{itemize}

\subsubsection*{Diagnóstico del frente bajo}

El par BB ($+0.83$) es \emph{positivo}: el modelo penaliza la formación
blanca. El par NB ($-2.80$) es el más favorecido.
Los coeficientes cúbicos pequeños ($c_{112}=-0.938$, $c_{122}=-0.127$)
implican que el dominio de los pares Go es casi completamente gobernado
por los términos lineales y cuadráticos, con poca modulación cúbica.
La ausencia de $H_1$ y robustez confirma que el polinomio no genera
complejidad topológica útil en el rango Go.

\subsubsection*{Energías}

\begin{center}
\begin{tabular}{ccccccc}
  \toprule
  Par & NN & N$\varnothing$ & NB & BN & B$\varnothing$ & BB & $\dE$ \\
  \midrule
  $H_{\text{F140}}$ & $-1.83$ & $-1.27$ & $-2.80$ & $+0.38$ & $+1.91$ & $+0.83$ & 4.72 \\
  \bottomrule
\end{tabular}
\end{center}

\newpage
% ==============================================================================
\section{Análisis comparativo: los 10 Hamiltonianos}
% ==============================================================================

\subsection{Tabla resumen de invariantes}

\begin{adjustbox}{max width=\textwidth}
\begin{tabular}{lllcccccr}
  \toprule
  \textbf{ID} & \textbf{Frente} & \textbf{Expresión (resumida)} &
  \textbf{$\mu$} & \textbf{$H_1$} & \textbf{$\dE$} &
  \textbf{Rob} & \textbf{Score} & \textbf{$c^*$ más cercano} \\
  \midrule
  $H_{\text{Alv}}$   & REF   & $xy$                     & 1 & 0.000 & 2.00 & 0.00 & 0.090 & 0.00 \\
  $H_{\text{M1}}$    & REF   & $x+2y-x^2y-xy^2$         & 2 & 0.000 & 2.00 & 0.33 & 0.190 & $\pm 1.24$ \\
  $H_{\text{F1b}}$   & 1     & {\small todos neg.\ + $-0.65x^2y$} & 2 & 0.120 & 11.75 & 1.00 & 0.417 & 0.42 \\
  $H_{\text{F1a}}$   & 1     & {\small $+0.85x^2y-0.80xy^2$}      & 2 & 0.166 & 8.00 & 1.00 & 0.454 & 0.06 \\
  $H_{\text{F2}}$    & 2     & {\small $+0.85x^2y+0.16xy^2$}      & 1 & 0.180 & 4.61 & 1.00 & 0.458 & 0.71 \\
  $H_{\text{F3}}$    & 3     & {\small $+2.50x-1.62y-2.78x^2$}    & 2 & 0.007 & 8.14 & 1.00 & 0.392 & $-0.47$ \\
  $H_{\text{F9}}$    & 9     & {\small $+1.997x+2.91y^2$}         & 2 & 0.000 & 11.02 & 0.00 & 0.090 & 1.19 \\
  $H_{\text{F25}}$   & 25    & {\small $+2.87x+2.83xy$}           & 2 & 0.000 & 11.73 & 0.00 & 0.090 & $-0.78$ \\
  $H_{\text{F70}}$   & 70    & {\small $+1.64x+2.15x^2$}          & 1 & 0.000 & 6.84 & 0.00 & 0.090 & $-0.27$ \\
  $H_{\text{F140}}$  & 140   & {\small $+1.59x-0.94x^2y$}         & 2 & 0.000 & 4.72 & 0.00 & 0.090 & $-1.24$ \\
  \bottomrule
\end{tabular}
\end{adjustbox}

\subsection{Jerarquía topológica vs.\ posición en el Hasse}

\begin{tcolorbox}[theorembox,title={Patrón general observado}]
Los Hamiltonianos del \textbf{Frente~1 y 2} comparten tres características
que los ubican en la cima del Hasse:
\begin{enumerate}[noitemsep]
  \item Los nodos $A_1$ están \emph{fuera} del rango de energía de los pares Go
        ($c^*$ mucho más negativo que $E_{\min}$), lo que garantiza que la
        fibra Go vive en la zona topológicamente más rica (entre el máximo
        local y el primer nodo).
  \item La separación $c^*_{\max} - E_{\min}$ es grande (8--12 unidades):
        el oval del dominio Go tiene larga vida en la filtración.
  \item El rango $\dE$ es grande (4--12) con buena uniformidad entre pares.
\end{enumerate}

Los Hamiltonianos del \textbf{Frente~9 al 140} fallan en:
\begin{itemize}[noitemsep]
  \item Un valor crítico $c^*$ cae \emph{dentro} del rango Go, cortando
        el ciclo topológico prematuramente ($H_1=0$).
  \item O bien la robustez es 0: pequeñas perturbaciones destruyen la
        estructura topológica.
\end{itemize}
\end{tcolorbox}

\subsection{Temperatura: resolución energética comparada}

La función de partición a temperatura $T=1$ (en unidades naturales) para
cada uno de los 10 Hamiltonianos es:

\begin{center}
\begin{tabular}{lrrrr}
  \toprule
  Hamiltoniano & $E_{\min}$ & $E_{\max}$ & $\dE$ & $T_{\text{Go}}\approx \dE/\ln 6$ \\
  \midrule
  $H_{\text{Alv}}$  & $-1.0$ & $+1.0$ & 2.0  & 1.12 \\
  $H_{\text{M1}}$   & $-1.0$ & $+1.0$ & 2.0  & 1.12 \\
  $H_{\text{F1b}}$  & $-12.76$ & $-1.02$ & 11.75 & 6.57 \\
  $H_{\text{F1a}}$  & $-9.63$ & $-1.63$ & 8.00  & 4.47 \\
  $H_{\text{F2}}$   & $-5.97$ & $-1.36$ & 4.61  & 2.58 \\
  $H_{\text{F3}}$   & $-7.95$ & $+0.19$ & 8.14  & 4.55 \\
  $H_{\text{F9}}$   & $-3.33$ & $+7.69$ & 11.02 & 6.16 \\
  $H_{\text{F25}}$  & $-6.83$ & $+4.90$ & 11.73 & 6.56 \\
  $H_{\text{F70}}$  & $-3.05$ & $+3.80$ & 6.84  & 3.82 \\
  $H_{\text{F140}}$ & $-2.80$ & $+1.91$ & 4.72  & 2.64 \\
  \bottomrule
\end{tabular}
\end{center}

$T_{\text{Go}} = \dE/\ln 6$ es la temperatura a la que la probabilidad del
par óptimo cae de $\approx 1$ a $\approx 1/2$: es la \textbf{temperatura
de decisión estratégica}. Los modelos del Frente~1 tienen la mayor
$T_{\text{Go}}$ de los candidatos con robustez~1, lo que significa que
discriminan estratégicamente bien incluso ante ruido térmico.

\subsection{Comparación visual de zonas críticas}

\begin{tikzpicture}[scale=1.0]
  % Eje de energía
  \draw[->,thick,cNavy] (-0.3,0) -- (13.5,0) node[right,cNavy] {$\dE$};
  \draw[thick,cNavy] (0,0.12) -- (0,-0.12) node[below,cNavy,font=\footnotesize] {0};

  % Cada Hamiltoniano como barra de rango energético
  \def\step{0.42}
  \foreach \ham/\label/\emin/\emax/\nA/\col/\i in {%
    HAlv/$H_\text{Alv}$/-1.0/1.0/1/cFrontRef/1,%
    HM1/$H_\text{M1}$/-1.0/1.0/2/cFrontRef/2,%
    HF1b/$H_\text{F1b}$/-12.76/-1.02/2/cFront1/3,%
    HF1a/$H_\text{F1a}$/-9.63/-1.63/2/cFront1/4,%
    HF2/$H_\text{F2}$/-5.97/-1.36/1/cFront2/5,%
    HF3/$H_\text{F3}$/-7.95/0.19/2/cFront2/6,%
    HF9/$H_\text{F9}$/-3.33/7.69/2/cFrontMid/7,%
    HF25/$H_\text{F25}$/-6.83/4.90/2/cFrontMid/8,%
    HF70/$H_\text{F70}$/-3.05/3.80/1/cFrontLow/9,%
    HF140/$H_\text{F140}$/-2.80/1.91/2/cFrontLow/10%
  }{%
    \pgfmathsetmacro{\ypos}{(11-\i)*\step}
    \pgfmathsetmacro{\xmin}{(\emin+13)/2}
    \pgfmathsetmacro{\xmax}{(\emax+13)/2}
    \filldraw[\col!70,draw=\col,line width=0.6pt,rounded corners=1pt]
      (\xmin,\ypos-0.1) rectangle (\xmax,\ypos+0.1);
    \node[left,font=\tiny,\col] at (\xmin-0.05,\ypos) {\label};
  }

  % Etiquetas del eje
  \foreach \val/\lbl in {0.5/{$-12$},2.0/{$-9$},3.5/{$-6$},5.0/{$-3$},6.5/{$0$},8.0/{$+3$},9.5/{$+6$}}{
    \draw[cGrey,thin] (\val,-0.05) -- (\val,0.05);
    \node[below,font=\tiny,cGrey] at (\val,0) {\lbl};
  }
  \node[below=0.35cm,cNavy,font=\small] at (6.5,0) {Energía en pares Go (unidades de $H$)};
\end{tikzpicture}

\medskip

Los Hamiltonianos del Frente~1 ($H_{\text{F1a}}$, $H_{\text{F1b}}$,
color verde) tienen sus barras energéticas completamente en el semiplano
negativo y son largas: el mínimo está muy por debajo de cero,
lo que genera la escala de temperatura estratégica más alta.

% ==============================================================================
\section{Zonas Go en \texorpdfstring{$\Gamma(H)$}{Gamma(H)}: síntesis}
% ==============================================================================

\subsection{Zonas estratégicas y su geometría}

Para cualquier Hamiltoniano de Go cúbico, las fibras de nivel
$H^{-1}(c)$ dividen el plano $(x,y)$ en regiones estratégicas:

\begin{tcolorbox}[defbox,title={Zonas estratégicas de $\Gamma(H)$}]
\begin{enumerate}[noitemsep]
  \item \textbf{Zona de piedras homogéneas} ($s_i=s_j$, NN o BB):
        pares en los vértices $(\pm 1,\pm 1)$.
        Si $H(\pm 1,\pm 1) < 0$: el modelo favorece la agrupación monocromática.
  \item \textbf{Zona de interacción mixta} ($s_i\neq s_j\neq 0$, NB o BN):
        pares en $(+1,-1)$ y $(-1,+1)$.
        Son los más frecuentes en Go real (frontera de territorios).
  \item \textbf{Zona de vacío} ($s_j=0$, N$\varnothing$ o B$\varnothing$):
        pares en $(\pm 1, 0)$.
        Si $H(s,0)=0$ (Alvarado): vacío invisible.
        Si $H(s,0)\neq 0$: el vacío es un elemento estratégico activo.
  \item \textbf{Punto crítico} (nodo $A_1$): punto de silla de $H$.
        Geometricamente, el paso de Go donde la estrategia cambia de
        ``local'' a ``global''; el mínimo y máximo local se separan.
\end{enumerate}
\end{tcolorbox}

\subsection{Estructura toroidal y ciclos de Vanishing}

El origen del toro en la fibra cúbica es la siguiente cadena de equivalencias:
\begin{align*}
  &\text{Curva cúbica plana genérica sobre }\C \;\cong\; \text{Toro }\C/\Lambda
  \quad (g=1, \;\text{fórmula de Riemann-Hurwitz}) \\
  &\text{Sobre }\R: \;\text{el toro real se manifiesta como oval compacto}
  + \text{rama no acotada} \\
  &\text{El oval }= \text{el ``ecuador'' del toro que contiene los pares Go}
\end{align*}
Cada nodo $A_1$ destruye uno de estos ovales: es el ``agujero del toro''
que se cierra. En términos Go: la zona de estrategia diferenciada
(el interior del oval) desaparece y el modelo ya no distingue entre
configuraciones de esa energía.

\subsection{Interpretación física: temperatura como exploración de la variedad}

\begin{tcolorbox}[theorembox,
  title={Temperatura como barrido de $\Gamma(H)$}]
A temperatura $T$, el sistema Go explora los pares de piedras con
distribución $p_{(s,s')}(T) \propto e^{-H(s,s')/T}$.

La \textbf{trayectoria de enfriamiento} $T:\infty\to 0$ corresponde,
en $\Gamma(H)$, a bajar desde el nivel $c=\langle H\rangle_{T=\infty}=\bar{E}$
(la media de los 6 pares) hasta $c=E_{\min}$.

\medskip

\textbf{Cada vez que esta trayectoria cruza un valor crítico $c^*$},
el sistema experimenta una transición de fase topológica:
cambia el número de ``modos estratégicos'' disponibles
(componentes conexas de la fibra de nivel).

\medskip

Los Hamiltonianos del \textbf{Frente~1} son superiores porque:
\begin{itemize}[noitemsep]
  \item La trayectoria de enfriamiento \emph{no cruza} ningún valor crítico
        (los nodos $A_1$ están por debajo de $E_{\min}$): el sistema
        baja suavemente al mínimo sin transiciones bruscas.
  \item La persistencia $H_1>0$ certifica que el oval topológico
        sobrevive todo el rango Go.
  \item La robustez $=1$ garantiza que esta propiedad no es un accidente
        de los coeficientes exactos.
\end{itemize}
\end{tcolorbox}

% ==============================================================================
\section{Conclusiones}
% ==============================================================================

\begin{enumerate}
  \item \textbf{Toda la geometría Go está en la fibra}.
        La variedad $\Gamma(H) = \{z=H(x,y)\}$ es trivialmente suave,
        pero la riqueza estratégica está en las fibras $H^{-1}(c)$:
        su topología cambia exactamente en los nodos $A_1$.

  \item \textbf{El Frente~1 maximiza el ovalo Go.}
        Los candidatos óptimos tienen sus nodos $A_1$ bien por debajo
        del mínimo energético de los pares Go, asegurando que el dominio
        Go esté contenido en la zona topológicamente más rica de la fibración
        de Milnor (el oval compacto, homeomorfo al ``ecuador del toro'').

  \item \textbf{$\mu=1$ puede superar a $\mu=2$ en $H_1$.}
        $H_{\text{F2}}$ con un solo nodo $A_1$ tiene la mayor persistencia
        ($H_1=0.180$), porque ese nodo único está perfectamente ubicado
        para maximizar la vida del ciclo Go.

  \item \textbf{La temperatura y el Milnor son duales.}
        El barrido de temperatura $T:\infty\to 0$ es equivalente, en
        $\Gamma(H)$, al barrido del parámetro $c:E_{\max}\to E_{\min}$.
        Cada cruce de $c^*$ es una transición de fase estratégica.
        El rango $\dE$ fija la escala $T_{\text{Go}}\approx \dE/\ln 6$.

  \item \textbf{El diagrama de Hasse ordena la calidad topológica.}
        Los frentes altos del Hasse corresponden a Hamiltonianos cuya
        fibración de Milnor es ``suave para Go'': los ciclos viven
        largo tiempo, la robustez es máxima, y la resolución térmica
        es grande. Los frentes bajos fallan en al menos uno de estos
        criterios.

  \item \textbf{El Hamiltoniano de referencia $H_{\text{M1}}$ es notable.}
        Con solo 2 nodos $A_1$ simétricos y $\dE=2$, la separación
        $\Delta c^*=2.48$ es significativa pero el rango Go es pequeño.
        El catálogo de 300 muestras encuentra mejores candidatos con
        mayor $\dE$ y mejor ubicación de nodos.
\end{enumerate}

\vspace{1cm}
\begin{center}
{\small\color{cGrey}
\textit{Pipeline}: \texttt{experiments/06\_hamiltonian\_families/} $\cdot$
\textit{Catálogo}: \texttt{output/catalog.json} $\cdot$
\textit{Diagrama de Hasse}: \texttt{output/figures/pareto\_fronts.json}
}
\end{center}

\end{document}
